\Needspace{18\baselineskip}
\section{Provenance and context closure at arbitrary box depth}
The shallow origin map alone is insufficient for deep replacement: an enclosing
box edge stays in place while occurrences in its body may change. We construct
provenance through each selected enclosing box and then induct on the typed site.

\Needspace{24\baselineskip}
\subsection{Local provenance for a changed label}
\begin{definition}[Fibre provenance]\label{term:fibre-provenance}
For old and new labels $l,l'$ at the same outer depth and a replacement $H$, a
fibre provenance certificate is an injective map
\[
 F_d(l')\longrightarrow F_d(l)+\operatorname{Occurrences}(H)
\]
preserving the hole observation in the indicated summand. It accounts for
changing one label while retaining or inserting the occurrences within that label.
\end{definition}
\lean{Graph/DeepContext}{12}{19}{Linear.FiberProvenance}
An unchanged label has the left inclusion as its origin map, which is injective
and preserves observation by reflexivity. If the new label merely has a proof
of equality to the old label, first transport its fibre address along that
equality and then use the left inclusion. Injectivity follows from transport
injectivity; the hole equation is the inverse direction of the transport
observation equation.
\lean{Graph/DeepContext}{22}{26}{Linear.FiberProvenance.unchanged}
\lean{Graph/DeepContext}{29}{33}{Linear.FiberProvenance.retained}

\begin{lemma}[Assembly of local provenance]
Suppose a signature-preserving per-edge relabelling of $G$ has a fibre provenance
certificate at every edge. If only a single designated edge can produce inserted
origins, these certificates assemble into provenance for the relabelled graph
relative to $G,H$.
\end{lemma}
\begin{proof}
At a result pair $(e,x)$ inspect its fibre origin. A retained origin $y$ gives
the old graph occurrence $(e,y)$; an inserted origin is already an occurrence
of $H$. Observation preservation is precisely the local certificate's equation.
For injectivity, consider the two output tags. Two retained origins identify
their outer edges and then their local retained addresses; local injectivity
identifies the original fibre addresses. Mixed output tags cannot agree.
If both origins are inserted, each outer edge must be the designated one:
otherwise the hypothesis that all its origins are retained contradicts its
inserted tag. With the outer edges identified, equality of inserted addresses
and local injectivity again identify the fibre addresses. Thus the global
origin map is injective in every case.
\end{proof}
\lean{Graph/DeepContext}{36}{91}{Linear.relabelProvenance}
The single-designated-edge restriction prevents two different outer edges from
claiming the same inserted occurrence. The lemma proves this separation; it
does not assume global injectivity as an extra graph law.

\Needspace{24\baselineskip}
\subsection{Lifting a body's provenance through one box}
Suppose the old and new bodies $B,B'$ have the same bound/free interface and a
provenance certificate $p$ for $B',B,H$. Their outer box labels differ only in
body content. Construct fibre provenance as follows. The shallow occurrence of
the box is retained as the old box's shallow occurrence. A new child address is
converted to the fibre-coordinate address of $B'$ and sent through $p$. Retained
body origins are converted back to body addresses and tagged as old box children;
inserted origins remain occurrences of $H$.
\begin{proof}[Injectivity and observation preservation]
A retained shallow box occurrence is distinct from every retained child and
every inserted origin by its tags. For two children, inspect the two body-origin
tags. Mixed tags cannot give equal lifted origins. If both are retained,
equality of the lifted children, followed by the old body's occurrence bijection,
identifies the retained body origins. If both are inserted, their equality is
already equality of body origins. In both cases body-provenance injectivity and
the new body's occurrence-bijection injectivity identify the original children.
For observations, the shallow box lookup is empty on both sides. A child's lookup
agrees with its body-coordinate observation by the bijection lemma; body
provenance gives the intermediate equality. In the retained case the inverse
bijection of the old body supplies the remaining observation equation. In the
inserted case the body equation already gives the replacement observation.
\end{proof}
\lean{Graph/DeepContext}{95}{147}{Linear.boxFiberProvenance}

To lift through a specified edge of a larger graph, use this box certificate at
that edge and an equality-transported retained certificate at every other edge.
Classical testing of edge equality selects the two cases. At the selected edge,
the given label equality identifies its old box, and the new label contains the
modified body with unchanged metadata. Elsewhere the label is unchanged, so all
local origins are retained. The assembly lemma therefore constructs provenance
for the existing concrete box-replacement operation.
\lean{Graph/DeepContext}{150}{170}{Linear.replaceBoxProvenance}

\Needspace{24\baselineskip}
\subsection{The deep context closure theorem}
\begin{theorem}[Provenance at every typed site]
For every typed site $s$ in $G$ and matching replacement $H$, the defined result
$\mathsf{plugAt}(s,H)$ has injective, hole-observation-preserving provenance
relative to $G,H$.
\end{theorem}
\begin{proof}
Induct on the site. A shallow site uses the concrete shallow provenance map.
For an inside site, the induction hypothesis gives provenance for the replacement
in the chosen body's inner site. Lift this certificate through the selected
outer box using the preceding construction. That is exactly the inside branch
of the definition of deep plugging. Induction follows the finite site address,
so the proof works at every finite nesting depth.
\end{proof}
\lean{Graph/DeepContext}{173}{180}{Linear.siteProvenance}

\begin{theorem}[Deep replacement preserves contexts under compatibility]
If $G,H$ have unique holes and the concrete site provenance is compatible, then
$\mathsf{plugAt}(s,H)$ has unique holes.
\end{theorem}
\begin{proof}
Use the constructed site provenance in the compatible-provenance uniqueness
theorem. No locality or occurrence-injectivity hypothesis remains to be supplied:
those properties have been proved for the actual replacement operation.
The remaining compatibility premise is explicitly a restriction on colliding
hole labels at retained origins.
\end{proof}
\lean{Graph/DeepContext}{183}{188}{Linear.plugAt_uniqueHoles_of_compatible}
Full disjointness of $G,H$ gives a simpler sufficient condition: instantiate the
full-disjointness provenance theorem with the same site certificate. This
condition is stronger than compatibility because it also inspects discarded host
occurrences; it is not asserted to be necessary.
\lean{Graph/DeepContext}{190}{194}{Linear.plugAt_uniqueHoles}
Every hole label of the deep result belongs to the inventory of the host or
replacement, by the provenance inventory-subset theorem. This is an upper bound
on the deep inventory, not an exact description of which host labels survive.
\lean{Graph/DeepContext}{197}{200}{Linear.plugAt_hasHole_subset}
If $H$ is hole-free, it has unique holes and is disjoint from $G$. Substituting
these two facts gives context closure for replacing any addressed edge by
hole-free material, with no separate freshness witness required.
\lean{Graph/DeepContext}{203}{207}{Linear.plugAt_holeFree_unique}
Finally, $\mathsf{Context.plugAt}$ accepts a certified host context, a typed site,
a certified replacement context, and full disjointness. It returns the concrete
deeply filled graph paired with its proved context certificate. The original
ordered outer interface and enclosing incidence laws are those already supplied
by deep replacement.
\lean{Graph/DeepContext}{209}{212}{Linear.Context.plugAt}

This provenance construction supplies globally unique-hole context closure and
concrete origin tracking through boxes. Section~\ref{sec:deep-inventory}
strengthens the inventory bound to an exact syntactic characterization and gives
a necessary-and-sufficient freshness condition at every finite box depth.
Associativity of plugging and recursive isomorphism of box contents remain
separate obligations. None of these results yet defines the concrete UAM.
